\documentclass[10pt,a4paper]{amsart}
\usepackage[margin=19mm]{geometry}
\usepackage{amsmath,amssymb,mathtools}
\usepackage[scr=rsfs,cal=euler]{mathalfa}
\usepackage{xfrac}
\usepackage[unicode,hidelinks,bookmarks=false]{hyperref}
\newcommand{\ie}{i.e.\ }
\newcommand{\cf}{cf.\ }
\newcommand{\resp}{respectively\ }
\newcommand{\Subsection}{\subsection}
\providecommand{\nobreakdash}{\nobreak\hbox{-}}
\setlength{\emergencystretch}{2em}
\multlinegap0pt
\usepackage{bm}
\usepackage{bbm}
\usepackage{dsfont}
\usepackage{xspace}
\usepackage{cancel}
\usepackage{mathtools}
\usepackage{amsthm}
\makeatletter
\@ifundefined{theorem}{\newtheorem{theorem}{Theorem}}{}
\makeatother
\newcommand{\bb}[1]{#1}
\title[Preconditioned normal equations for solving discretised part]{Preconditioned normal equations for solving discretised partial differential equations}
\author{Lorenzo Lazzarino, Yuji Nakatsukasa, Umberto Zerbinati}
\date{Source: arXiv:2502.17626v3 (CC BY 4.0)}
\begin{document}
\maketitle

\noindent\emph{Source and licence.} Preconditioned normal equations for solving discretised partial differential equations. Version arXiv:2502.17626v3 (CC BY 4.0), distributed under the Creative Commons Attribution 4.0 International licence, as recorded in the arXiv licence metadata for this exact version and shown on the abstract page https://arxiv.org/abs/2502.17626v3. Full rights notice: licenses/arxiv-2502.17626-CC-BY-4.0.txt.\par\smallskip

\noindent\textit{Editorial scope.} Counted unit: the Theorem whose source label is \texttt{thm:conv} in the article (source0.tex lines 420--432, source label \texttt{thm:conv}), delivered together with the article's own complete original proof of it (source0.tex lines 434--477, one \texttt{proof} environment of 44 lines). No mathematics is written, repaired or re-derived: statement and proof are the article's own text line for line. Required context retained uncounted: eq:CGbound (lines 187--190). Every definition, display and label of those lines is retained unchanged. Omitted from this package and not redistributed: 1--186, 191--419, 433--433, 478--986. Nothing inside a retained excerpt is omitted or altered: every retained body is delivered exactly as the article's lines are written, so no omission receipt of any kind accompanies this package. Citations: the retained excerpts carry no citation command at all, so no bibliography entry of the article is transcribed and no reference is redistributed. Reference adaptation: none was needed and none was made; every reference inside the delivered unit resolves inside it, so no unresolved reference or citation is produced. Printed slips kept literal: no printed word, symbol, hypothesis or number of the counted statement or of its proof was altered. Layout and class: the article's own class is \texttt{siamart220329} with packages including mathdots, todonotes, booktabs, subcaption, multirow, lipsum, amsfonts, graphicx, epstopdf, mathtools, algorithmic, ulem, tikz, pgfplots; the wrapper is the lane's pinned \texttt{amsart} editorial wrapper at 10pt on A4, which restates the theorem-like environments the retained text needs and adds the article's own macro definitions that the retained text needs (\texttt{\textbackslash bb}), with extra wrapper packages (bm, bbm, dsfont, xspace, cancel, mathtools). The delivered numbering is the unit's own. Assets: no retained span contains a figure environment, a graphics-inclusion command, a PDF-page inclusion, a table or a TikZ picture; no image file is copied into this package, the package declares no asset dependency and no image file is redistributed with it. Licence: version arXiv:2502.17626v3 of this article is distributed under the Creative Commons Attribution 4.0 International licence, as recorded in the arXiv licence metadata for this exact version and shown on the abstract page https://arxiv.org/abs/2502.17626v3. A full rights notice is delivered at licenses/arxiv-2502.17626-CC-BY-4.0.txt. Characters: every retained excerpt is pure ASCII, so the delivered engine, XeLaTeX with its default Latin Modern text font, reports no missing character.
\begin{equation}
    \label{eq:CGbound}
    \frac{\|x^*-x_k\|_B}{\|x^*-x_0\|_B} \leq \underset{p\in P_k}{\bb{\min}} \underset{\lambda \in \Lambda(B)}{\bb{\min}} |p(\lambda)| \leq 2 \left(\frac{\sqrt{\kappa(B)}-1}{\sqrt{\kappa(B)}+1} \right)^k = 2 \left(\frac{\kappa(A)-1}{\kappa(A)+1} \right)^k,
\end{equation}

\begin{theorem}
    \label{thm:conv}
    The relative error at the $k$-th iteration of the CG method applied to the preconditioned equation
    \begin{equation}
        \label{eq:PrecnormalPDE}
        G_T^{-1}B_T\vec{x} = G_T^{-1}A^TT\vec{b}, \quad G_T = P^TTP, \quad B_T = A^TTA,
    \end{equation}
    with exact solution $x^*$, is such that
    \begin{equation}
        \frac{\|x^* - x_k \|_{B_T}}{\|x^* - x_0 \|_{B_T}} \leq \underset{p\in P_k}{\bb{\min} \underset{\,\,\sigma \in \Sigma_T(AP^{-1})}{\bb{\max}}} |p(\sigma^2)| \leq 2 \left( \frac{\kappa_T(AP^{-1}) -1}{\kappa_T(AP^{-1}) +1} \right)^k,
    \end{equation}
    where $P_k$ is the space of polynomials $p$ of degree less than or equal to $k$ and unit lowest order coefficient.
\end{theorem}

\begin{proof}
    Recall that since $T$ is symmetric positive definite, it always exists the matrix $T^{1/2}$, that is as well symmetric positive definite. Observe that, \bb{left multiplying by}  $T^{1/2}P$, we can rewrite equation~\eqref{eq:PrecnormalPDE} as:
    \begin{equation}
        \label{eq:PrecnormalPDE-rewritten}
        T^{-1/2}P^{-T}A^T T A P^{-1} T^{-1/2} y = T^{-1/2}P^{-T}A^T Tb, \quad x=P^{-1}T^{-1/2}y.
    \end{equation}
    To simplify the notation, let us define the matrix $Z := P^{-T}A^T T A P^{-1}$.
    Further, denoting by $y^*$ the exact solution of equation~\eqref{eq:PrecnormalPDE-rewritten}, the following relation holds
    \begin{equation}
        \begin{aligned}
            \|y^*- y\|_{T^{-1/2}Z T^{-1/2}} & = (y^*-y)^TT^{-1/2}Z T^{-1/2}(y^*-y)                         \\
                                            & = (x^*-x)^T(P^T T^{1/2}) T^{-1/2}Z T^{-1/2}(T^{1/2}P)(x-x^*) \\
                                            & = (x^*-x)^T P^T Z P(x-x^*)                                   \\
                                            & = (x^*-x)^T A^TTA(x-x^*)                                     \\
                                            & = \|x^*-x\|_{A^TTA}.
        \end{aligned}
    \end{equation}
    Moreover, by the definition of $T$-singular values given above, and by the similarity transformation
    \begin{equation}
        T^{-1}Z = T^{-1/2}(T^{-1/2}ZT^{-1/2})T^{1/2},
    \end{equation}
    we have that
    \begin{equation}
        \label{eq:singVsEige}
        %\sigma^2_{T,i}\left(AP^{-1}\right) = \lambda_i\left((AP^{-1})^\circ (AP^{-1}) \right) = \lambda_i\left(T^{-1}Z\right) = \lambda_i \left( T^{-1/2}ZT^{-1/2} \right).
        \sigma^2_{T,i}\left(AP^{-1}\right) = \lambda_i\left(T^{-1}Z\right) = \lambda_i \left( T^{-1/2}ZT^{-1/2} \right).
    \end{equation}
    Thus,
    \begin{equation}
        \begin{aligned}
            \frac{\|x^* - x_k \|_{B_T}}{\|x^* - x_0 \|_{B_T}} & = \frac{\|y^* - y_k\|_{T^{-1/2}ZT^{-1/2}}}{\|y^* - y_0\|_{T^{-1/2}ZT^{-1/2}}}                                     \\
                                                              & \leq \underset{p\in P_k}{\bb{\min}} \underset{\,\,\lambda \in \Lambda(T^{-1/2}ZT^{-1/2})}{\bb{\min}} |p(\lambda)| \\
                                                              & = \underset{p\in P_k}{\bb{\min}} \underset{\,\,\sigma \in \Sigma_T(AP^{-1})}{\bb{\min}} |p(\sigma^2)|
        \end{aligned}
    \end{equation}

    Finally, combining \eqref{eq:CGbound} and \eqref{eq:singVsEige}, we obtain
    \begin{equation}
        \begin{aligned}
            \frac{\|x^* - x_k \|_{B_T}}{\|x^* - x_0 \|_{B_T}} & \leq 2\left( \frac{\sqrt{\kappa(T^{-1/2}ZT^{-1/2})} -1}{\sqrt{\kappa(T^{-1/2}ZT^{-1/2})} +1} \right)^k \\
                                                              & = 2\left( \frac{\sqrt{\kappa(T^{-1}Z)} -1}{\sqrt{\kappa(T^{-1}Z)} +1} \right)^k
        \end{aligned}
    \end{equation}
\end{proof}

\end{document}
